\begin{frame}{SOAP frente a REST}
\small
\begin{tabularx}{\linewidth}{@{}l X X@{}}
\toprule
 & SOAP & REST \\
\midrule
Contrato & WSDL obligatorio \cite{w3c_wsdl11} & OpenAPI opcional \cite{openapi31} \\
Formato & XML dentro de un sobre & JSON casi siempre \\
Errores & SOAP Fault, enviado con HTTP 500 & Códigos HTTP: 404, 409, 503\dots{} \cite{rfc9110} \\
Caché & Casi nunca: las peticiones van por \texttt{POST} & Sí: \texttt{GET} con \texttt{Cache-Control} y \texttt{ETag} \cite{rfc9111} \\
Tamaño en nuestro PFC & 703 bytes de respuesta y 308 de petición & 252 bytes de respuesta \\
Sistemas reales & SRI, facturación electrónica & DynamoDB y nuestros tres proyectos \\
\bottomrule
\end{tabularx}
\end{frame}

\begin{frame}{Postura del equipo}
\begin{columns}[T]
\begin{column}{0.47\textwidth}
\begin{block}{SOAP conviene cuando\dots}
\begin{itemize}
\item muchas organizaciones deben cumplir un contrato fijo, como los emisores de facturas ante el SRI;
\item el documento ya es XML.
\end{itemize}
\end{block}
\end{column}
\begin{column}{0.47\textwidth}
\begin{block}{REST conviene cuando\dots}
\begin{itemize}
\item el cliente es una web o una app móvil;
\item importa el tamaño del mensaje y aprovechar la caché de HTTP.
\end{itemize}
\end{block}
\end{column}
\end{columns}
\vspace{10pt}
\begin{alertblock}{Pero REST mal hecho cuesta}
Solo el 0,8\,\% de 500 API públicas cumple todos los principios REST \cite{neumann2021}. Romper las reglas de diseño hace la API más difícil de entender en 11 de 12 reglas probadas \cite{bogner2023}, y cambiarla sin romper a los clientes exige versiones y compatibilidad \cite{lercher2024}.
\end{alertblock}
\end{frame}
